\documentclass{article}
\usepackage[english]{babel}
\usepackage[a4paper,top=2cm,bottom=2cm,left=3cm,right=3cm,marginparwidth=1.75cm]{geometry}
\usepackage{amsmath}
\usepackage{graphicx}
\usepackage[hidelinks]{hyperref}
\usepackage{listings}
\usepackage{xcolor}
\lstset{
  language=SQL,
  basicstyle=\ttfamily,
  keywordstyle=\color{blue}\bfseries,
  stringstyle=\color{orange},
  commentstyle=\color{green!50!black},
  showstringspaces=false,
  numbers=left,
  numberstyle=\tiny,
  stepnumber=1,
  numbersep=8pt
}

\title{New-Horizons Report}

\author{Nicon-David Milandru}

\begin{document}

\maketitle

\tableofcontents

\section{Introduction}
Travel agencies face complexity in managing \textbf{packages}, \textbf{bookings}, and \textbf{client information} due to multiple interrelated entities.

This project presents the design and implementation of a relational \textbf{database system} for a travel agency. The purpose of the system is to manage the entire lifecycle of \textbf{clients}, \textbf{bookings}, \textbf{invoices}, and the structure of \textbf{travel packages}. The scope of the project covers conceptual modeling, relational schema generation, sample data population, and validation queries.

The \textbf{database} was modeled using \textbf{Oracle SQL Developer Data Modeler}, implemented in \textbf{SQL}, and tested via \textbf{Supabase}. Limitations include the absence of a front-end application and performance optimization for very large datasets.

\section{Logical Design}

\subsection{Entities}
The \textbf{entity-relationship diagram (ERD)} was created using \textbf{Oracle SQL Data Modeler} and includes relevant \textbf{entities} for the project. Each \textbf{entity} contains well-defined \textbf{attributes} with appropriate \textbf{data types}, which may be optional or mandatory, and may also serve as \textbf{primary keys} or \textbf{foreign keys}. The \textbf{schema} is fully normalized. Many-to-many relationships between \textbf{entities} were resolved through the introduction of \textbf{associative entities}.

\subsection{Relationships}
The \textbf{Client} entity has a mandatory one-to-many relationship with the \textbf{Booking} entity (a client must have one or more \textbf{bookings}).

The \textbf{Booking} entity has a mandatory many-to-one relationship with the \textbf{Client} entity (each \textbf{booking} is made by exactly one \textbf{client}). It also has a mandatory one-to-many relationship with the \textbf{Invoice} entity (each \textbf{booking} is associated with one or more \textbf{invoices}). Finally, it has a mandatory many-to-one relationship with the \textbf{Travel Package} entity (each \textbf{booking} contains exactly one \textbf{travel package}).

The \textbf{Invoice} entity has a mandatory many-to-one relationship with the \textbf{Booking} entity (each \textbf{invoice} is associated with exactly one \textbf{booking}).

The \textbf{Travel Package} entity has an optional one-to-many relationship with the \textbf{Booking} entity (a \textbf{travel package} may appear in zero or more \textbf{bookings}). It also has one-to-many relationships with the \textbf{associative entities} \textbf{Accommodation Package}, \textbf{Destination Package}, \textbf{Transport Package}, and \textbf{Group Package}, which resolve potential many-to-many relationships. The relationships with \textbf{Accommodation Package}, \textbf{Destination Package}, and \textbf{Transport Package} are mandatory (each \textbf{travel package} must include at least one \textbf{destination}, \textbf{accommodation}, and \textbf{transport}). However, the relationship with \textbf{Group Package} is optional, as participation in a \textbf{group activity} is not required.

The \textbf{Destination Package} entity has a mandatory many-to-one relationship with the \textbf{Travel Package} entity and a mandatory many-to-one relationship with the \textbf{Destination} entity. The \textbf{Destination} entity has an optional one-to-many relationship with the \textbf{Destination Package} entity.

The same logic applies to the \textbf{Accommodation Package} and \textbf{Transport Package} entities. Each of these \textbf{associative entities} has a mandatory many-to-one relationship with the \textbf{Travel Package} entity and a mandatory many-to-one relationship with its respective \textbf{entity} (\textbf{Accommodation} or \textbf{Transport}). Similarly, the \textbf{Accommodation} and \textbf{Transport} entities have optional one-to-many relationships with their corresponding \textbf{associative entities}.

These relationships, together with the introduction of \textbf{associative entities}, resolve the potential many-to-many relationships, which are not considered good practice in a \textbf{logical relational model}.

The \textbf{Group Package} entity has a mandatory many-to-one relationship with the \textbf{Travel Package} entity and an optional one-to-many relationship with the \textbf{Guide Package} entity. This entity introduces details about \textbf{group activities} and also an \textbf{associative entity} to represent the assignment of a \textbf{guide} to a \textbf{group activity}.

The \textbf{Guide Package} entity serves as the \textbf{associative entity} between \textbf{Group Package} and \textbf{Guide}. It has a mandatory many-to-one relationship with the \textbf{Group Package} entity and a mandatory many-to-one relationship with the \textbf{Guide} entity. The \textbf{Guide} entity has an optional one-to-many relationship with the \textbf{Guide Package} entity. This structure prevents a direct many-to-many relationship between \textbf{guides} and \textbf{group activities}.

\section{Relational Design}
The \textbf{relational model} was automatically generated from the \textbf{logical model (ERD)} using the \textbf{Engineer} function provided by \textbf{Oracle SQL Data Modeler}. The resulting \textbf{schema} provides additional details about the \textbf{database structure}, including \textbf{attribute data types}, \textbf{primary keys}, and \textbf{inherited keys}.

\section{Implementation}

\subsection{Data Definition Language}
From the \textbf{relational model}, \textbf{Oracle SQL Data Modeler} allows the generation of a \textbf{Data Definition Language (DDL)} script. This is \textbf{SQL} code that creates \textbf{tables} for all \textbf{entities}, \textbf{attributes}, and \textbf{relationships} defined in the \textbf{schema}. For example, the code generated for the \textbf{Client} entity is:

\begin{lstlisting}
CREATE TABLE CLIENT (
    client_id INTEGER GENERATED ALWAYS AS IDENTITY PRIMARY KEY,
    first_name VARCHAR(50) NOT NULL,
    last_name VARCHAR(50) NOT NULL,
    email VARCHAR(100) NOT NULL UNIQUE,
    phone_number VARCHAR(20),
    date_of_birth DATE,
    passport_number VARCHAR(20),
    registration_date DATE NOT NULL,
    loyalty_points INTEGER NOT NULL
);
\end{lstlisting}

\subsection{Sample Data}

After generating the \textbf{DDL}, I deployed the \textbf{database} on \textbf{Supabase} in order to execute my scripts. Creating a \textbf{project} and running \textbf{SQL scripts} on the platform is straightforward and intuitive. Once the \textbf{DDL} was executed, I prepared \textbf{sample data} to populate the \textbf{tables}. For instance, the sample data for the \textbf{Destination} entity is as follows:

\begin{lstlisting}
-- DESTINATIONS
INSERT INTO DESTINATION (country, city, description, airport_code)
VALUES
('France', 'Paris', 'City of Light', 'CDG'),
('Italy', 'Rome', 'The Eternal City', 'FCO'),
('Japan', 'Tokyo', 'Modern metropolis', 'HND');    
\end{lstlisting}

\subsection{Queries}
With the \textbf{sample data} loaded into the \textbf{database}, we can now query information and test the \textbf{structure}. I created a number of \textbf{SQL queries} that retrieve meaningful data from the \textbf{tables}. For example, the following query calculates the average \textbf{guide rating} by \textbf{travel package}:

\begin{lstlisting}
SELECT 
    tp.package_name,
    AVG(g.rating) AS avg_guide_rating
FROM travel_package tp
JOIN group_package gp ON tp.package_id = gp.package_id
JOIN guide_package gpk ON gp.group_package_id = gpk.group_package_id
JOIN guide g ON gpk.guide_id = g.guide_id
GROUP BY tp.package_name
ORDER BY avg_guide_rating DESC;    
\end{lstlisting}

\section{Reflection}
This is a simple project that also has real-life applications, involving multiple tools, and developing a variety of skills. 

Starting from the design of a \textbf{logical model}, we must consider how to adapt it to a \textbf{real-world scenario}, carefully think about \textbf{relationships}, and determine the \textbf{information} we want to store in the \textbf{database}. Transitioning to a \textbf{relational model}, we analyze how the \textbf{information} interacts with itself (\textbf{relationships}), how to distinguish it (\textbf{keys}), and how to store it efficiently (\textbf{data types}). Next, we learn to use \textbf{tools} that generate \textbf{SQL scripts} from a \textbf{schema}, which forms the \textbf{implementation phase}. Finally, we use the \textbf{model} and the entered \textbf{data} to retrieve \textbf{information} through \textbf{queries}. 

Ultimately, this project is about \textbf{information}: how it flows, what it represents, how it is stored, how it interacts, and how we can extract it. It is essentially the process of deriving \textbf{information} from \textbf{information}.
\end{document}